\section{The Model}\label{model}

We start with a generalised form of the Lotka-Volterra equation 
\begin{equation}\label{lotka-volterra}
\dot{\bn} = \br*\bn + \bn*\bbeta\bn.
\end{equation}
Here \bn\ is the population density, the component $n_i$ being the
number of individuals of species $i$, \br\ is the difference
between reproduction and death, \bbeta\ is the interaction matrix,
with $\beta_{ij}$ being the interaction between species $i$ and $j$, *
referring to elementwise multiplication and {\tt mutate} is the
mutation operator.

\subsection{Lotka-Volterra Dynamics}

The most obvious thing about equation (\ref{lotka-volterra}) is its
fixed point 
\begin{equation}\label{fixed point}
\hat{\bn} = -\bbeta^{-1}\br,
\end{equation}
where $\dot{\bn}=0$. For this point to be biologically meaningful, all
components of $\hat{\bn}$ must be positive, giving rise to the following
inequalities:
\begin{equation}\label{positive species}
\hat n_i = \left(\bbeta^{-1}\br\right)_i>0, \forall i
\end{equation}
The stability of this point is related to the
negative definiteness of derivative of $\dot{\bn}$ at $\hat{\bn}$. The
components of the derivative are given by
\begin{equation}\label{derivative}
\frac{\partial\dot{n}_i}{\partial n_j} =
\delta_{ij}\left(r_i+\sum_k\beta_{ik}n_k\right) + \beta_{ij}n_i
\end{equation}
Substituting eq (\ref{fixed point}) gives
\begin{equation}
\left.\frac{\partial\dot{n}_i}{\partial n_j}\right|_{\hat{\bn}}=
-\beta_{ij}\left(\bbeta^{-1}\br\right)_i
\end{equation}

Stability of the fixed point requires that this matrix should be
negative definite. Since the $\left(\bbeta^{-1}\br\right)_i$ are
all negative by virtue of (\ref{positive species}), each minor
determinant of this matrix is equal to a minor determinant of \bbeta\
multiplied by a positive number, stability of the equilibrium is
equivalent to \bbeta\ being negative definite.

A weaker condition is to require that the system remain bounded with
time:
\begin{equation}\label{boundedness}
\sum_i\dot{n_i}=\br\cdot\bn + \bn\cdot\bbeta\bn < 0, \;\forall \bn:
\sum_in_i>N \;\exists N
\end{equation}

As \bn\ becomes large in any direction, this functional is dominated
by the quadratic term, so this implies that  $\bn\cdot\bbeta\bn\leq0
\; \forall\bn: n_i>0$. Negative definiteness of \bbeta\ is sufficient,
but not necessary for this condition. For example, the predator-prey
relations (heavily normalised) have the following matrix as \bbeta:
\begin{math}
\bbeta=\left(\begin{array}{cc}
-1 & 2\\
-2 & 0\\
\end{array}\right)
\end{math}
which has eigenvalues $3/2, -5/2$. If we let $\bn=(x,y), x,y\geq0$, then
$\bn\cdot\bbeta\bn=-2x^2$, which is clearly non-positive
for all nonpositive $x$.

Consider adding a new row and column to \bbeta. What is condition is
the new row and column required to satisfy such that equation
(\ref{boundedness}) is satisfied. Break up \bbeta\ in the following
way:
\begin{displaymath}
\left(
  \mbox{
     \begin{tabular}{c|c}
       $\begin{array}{ccc}\ddots\\&{\bf A}\\&&\ddots\end{array}$ & 
       $\begin{array}{c}\vdots\\{\bf B}\\\vdots\end{array}$ \\
       \hline
       $\begin{array}{ccc}\cdots&{\bf C}&\cdots\end{array}$ & D
     \end{tabular}
   }
\right)
\left(
  \mbox{
     \begin{tabular}{c}
     $\begin{array}{c}\vdots\\{\bf n_1}\\\vdots\end{array}$\\
     \hline
     $n_2$
     \end{tabular}
    }
\right)
\end{displaymath}

Condition (\ref{boundedness}) becomes:
\begin{equation}\label{boundedness2}
{\bf n_1}\cdot{\bf A}{\bf n_1} + {\bf n_1}\cdot({\bf B}+{\bf C})n_2 +
Dn_2^2 \leq 0
\end{equation}

Let 
\begin{displaymath}
a=\max_{n=1} \bn\cdot A\bn,\mbox{ and } b=\max_{i}B_i+C_i.
\end{displaymath}
 Then a sufficient but
not necessary condition for condition (\ref{boundedness2}) is
\begin{displaymath}
an_1^2+bn_1n_2+Dn_2^2\leq0
\end{displaymath}

The maximum value with respect to $n_2$ is $an_1^2-(bn_1)^2/4D$, so
this requires that
\begin{equation}\label{boundedness3}
b < 2\sqrt{ad}
\end{equation}

\subsection{Mutation}

With mutation, equation (\ref{lotka-volterra}) reads
\begin{equation}
\dot{\bn} = \br*\bn + \bn*\bbeta\bn + {\tt mutate}(\bmu,\br,\bn).
\end{equation}


The difficulty with adding mutation to this model is how to define the
mapping between genotype space and phenotype space, or in other words,
what defines the {\em embryology}. A few studies, including Ray's
Tierra world, do this with an explicit mapping from the genotype to to
some particular organism property (e.g. interpreted as machine language
instructions, or as weight in a neural net). These organisms then
interact with one another to determine the population dynamics. In
this model, however, we are doing away with the organismal layer, and
so an explicit embryology is impossible. The only possibility left is
to use a statistical model of embryology. The mapping between
genotype space and the population parameters \br,
\bbeta{}  is expected to look like a rugged
landscape, however, if two genotypes are close together (in a Hamming
sense) then one might expect that the phenotypes are likely to be
similar, as would the population parameters. This I call {\em random
embryology with locality}.

  In the simple case of point mutations, the probability $P(x)$ of any
child lying distance $x$ in genotype space from its parent follows a
Poisson distribution. Random embryology with locality implies that the
phenotypic parameters are distributed randomly about the parent
species, with a standard deviation that depends monotonically on the
genotypic displacement. The simplest such model is to distribute the
phenotypic parameters in a Gaussian fashion about the parent's values,
with standard deviation proportional to the genotypic
displacement. This constant of proportionality can be conflated with
the species' intrinsic mutation rate, to give rise another phenotypic
parameter \bmu.  It is assumed that the probability of a mutation generating
a previously existing species is negligible, and can be ignored. We
also need another arbitrary parameter $\rho$, ``species radius'',
which can be understood as the minimum genotypic distance separating
species, conflated with the same constant of proportionality as \bmu.

In summary, the mutation algorithm is as follows:
\begin{enumerate}
\item The number of mutant species arising from species $i$ within a
timestep is $\mu_i\alpha_in_i/\rho$. This number is rounded
stochastically to the nearest integer, e.g. 0.25 is rounded up to 1
25\% of the time and down to 0 75\% of the time.

\item Roll a random number from a Poisson distribution
$e^{-x/\mu+\rho}$ to determine the standard deviation $\sigma$ of phenotypic
variation. 

\item Vary \br\ according to a Gaussian distribution about the
parents' values, with $\sigma\alpha_0$ as the standard deviation,
where $\alpha_0$ is the range of values that \br\ is initialised
to, ie $\alpha_0$=\verb|repro_rate(random,maxval)|$-$\verb|repro_rate(random,minval)|

\item The diagonal part of \bbeta\ must be negative, so vary \bbeta\
according to a log-normal distribution. This means that if the old
value is $\beta$, the new value becomes
$\beta'=-\exp(\log_e(\beta)+\sigma)$. These values cannot become
arbitrarily small, however, as this would imply that some species make
arbitrarily small demands on the environment, and will become infinite
in number. In ecolab, the diagonal interactions terms prevented from
becoming larger than $-r/(.1*{\tt INT\_MAX})$.

\item The offdiagonal components of \bbeta, are varied in a similar
fashion to \br. However new connections are added, or old ones
removed according to $\lfloor 1/r\rfloor$, where $r\in(-2,2)$ is
chosen from a uniform distribution. The values on the new connections
are chosen from the same initial distribution that the offdiagonal
values where originally set with, ie the range
\verb|interaction(offdiag,random,minval)| to
\verb|interaction(offdiag,random,maxval)|. Since $a$ in condition
(\ref{boundedness3}) is computationally expensive, we use a slightly stronger
criterion that is sufficient, computationally tractable yet still
allows ``interesting'' non-definite matrix behaviour namely that the sum
$\beta_{ij}+\beta_{ji}$ should be nonpositive.


\item \bmu\ must be positive, so should evolve according to the
log-normal distribution like the diagonal components of
\bbeta. Similar to \bbeta, it is a catastrophe to allow \bmu\ to
become arbitrarily large. In the real world, mutation normally exists
at some fixed background rate --- species can reduce the level of
mutation by improving their genetic repair algorithms. In ecolab, this
ceiling on \bmu\ is given by the \verb|mutation(random,maxval)| variable.

\end{enumerate}

\section{Structure of the Ecolab Simulation System}

Ecolab is built on three main class libraries: the {\tt arrays} class,
which implements dynamic arrays, the {\tt globals} class, which define
the global state variables of the model and {\tt tcl++} which provides
bindings to the TCL scripting language. Each of these class libraries
will be described in subsequent sections.

Further modules of the system include {\tt analysis} and {\tt
  Xecolab}, which provides a number of generic instruments for
observing the Ecolab system, and {\tt aux}, which contains code for
  implementing  a
client/server version of Ecolab, with the
instruments running in the client system, and the model running on the
server, and code for implementing checkpoint/restart. This minimises the amount of traffic between a compute server
and a client.

The Ecolab model itself is defined in the model specific file {\tt
  ecolab.cc}. To define another model, replace this file by another
  one with similar functions. An example of this is {\tt
  newman.cc}. See \S\ref{new model} for more details.

The whole computation is constructed from a TCL\cite{Ousterhout94}
script. Example scripts include {\tt ecolab.tcl}, {\tt pred-prey.tcl},
{\tt lifetime.tcl} and {\tt console.tcl/engine.tcl} for a sample
client/server system.

\begin{figure}
\input{pst-node}
\pspicture(2,0)(15,7)
\rput(7,7){\rnode{ecolabtcl}{\psframebox{ecolab.tcl}}}
\rput(5,5){\rnode{modeltcl}{\psframebox{model.tcl}}}
\rput(7,5){\rnode{Xecolab}{\psframebox{Xecolab.tcl}}}
\rput(9,5){\rnode{ecolabcc}{\psframebox{ecolab.cc}}}
\rput(5,3){\rnode{analysis}{\psframebox{analysis.cc}}}
\rput(12.5,3){\rnode{globals}{\psframebox{globals.cc}}}
\rput(10,3){\rnode{initarrays}{\psframebox{init\_arrays.cc}}}
\rput(5,1){\rnode{tclpp}{\psframebox{tcl++.h}}}
\rput(7,1){\rnode{arrays}{\psframebox{arrays.h}}}
\rput(7,0){\rnode{carrays}{\psframebox{c\_arrays.c}}}
\ncline{ecolabtcl}{modeltcl}
\ncline{ecolabtcl}{Xecolab}
\ncline{ecolabtcl}{ecolabcc}
\ncline{Xecolab}{analysis}
\ncline{ecolabcc}{tclpp}
\ncline{ecolabcc}{globals}
\ncline{ecolabcc}{arrays}
\ncline{ecolabcc}{initarrays}
\ncline{analysis}{tclpp}
\ncline{globals}{arrays}
\ncline{initarrays}{arrays}
\ncline{arrays}{carrays}
\endpspicture
\caption{Structure of the Ecolab Simulation System}
\end{figure}

